% Introduction to Part I. Unnumbered text: no \label except the \phantomsection target.
\phantomsection\label{intro:I}
Part~I studies the linear operators of Ba\l aban's renormalization group at the trivial background $U=1$, on a single fine torus $\Tf$ of spacing $\eta=L^{-k}$ and side $n$, with averages over the unit blocks $B(y)$, $y\in\T$ (Section~\ref{sec:setup}). At $U=1$ every operator acts componentwise in a basis of the Lie algebra, so the fields are real-valued, and every constant of Part~I depends only on the coefficient $a$ of the averaging term and on the H\"older exponent $\alpha$, not on $L$, $k$ or $n$ (Section~\ref{sec:intro}). Parts~II and~III write every operator at a background as its counterpart at $U=1$ plus a perturbation (Section~\ref{sec:perturbation}), so Part~I is the base of the whole operator theory.

\paragraph{What this Part proves.} Bounds uniform in $k$: coercivity, exponential decay between unit blocks, and, for the Green functions, the minimizer maps and the constrained Green functions, pointwise bounds on values, first derivatives and H\"older quotients. They hold for the scalar Green functions $G_\sigma=(-\Lap+aQ_\sigma^*Q_\sigma)^{-1}$ (Theorem~\ref{thm:scalar-green}); for the averaged operators $(Q_\sigma G_\sigma Q_\sigma^*)^{-1}$, the minimizer maps, the effective actions, the constrained Green functions and the projection $R$ of Ba\l aban's Landau-type gauge (Section~\ref{sec:scalar-avg}); for the gauge-field operator $\bL_a=d^*d+dRd^*+a\bQ^*\bQ$ and its inverse $\bG$ (Theorems~\ref{thm:gauge-coercive} and~\ref{thm:gauge-green}); and for the gauge-field minimizer map $\bH$, effective action $\bD$ and constrained Green function $\bC$ (Section~\ref{sec:gauge-min}). Theorems~\ref{thm:A}--\ref{thm:E} of the Introduction collect these results. The tools are a block Poincar\'e inequality (Section~\ref{sec:poincare}), one Fourier lower bound (Corollary~\ref{cor:QKQ-lower}), the Combes--Thomas method (Section~\ref{sec:decay}) and a lattice Schauder estimate (Theorem~\ref{thm:local-regularity}). The coercivity of $\bL_a$ is proved in position space; its key step is Lemma~\ref{lem:Q-divergence}, by which $\bQ$ nearly separates exact and co-closed fields.

\paragraph{Imports.} Appendix~\ref{app:green} gives the decay of the Green function of $\Z^d$ and of its differences, hence that of $\eta\Z^d$ (Proposition~\ref{prop:green-eta}). The averages $\bQ$ and the projection $R$ are those of \cite{BalPropI}. No proof in Part~I uses a later result. Section~\ref{sec:setup} points ahead only to say where $\bQ$ comes from: it is the linearization at $U=1$ of the composite of $k$ gauge-adapted one-step averagings (Proposition~\ref{prop:linearized-averages}(c)).

\paragraph{Exports.}
\begin{itemize}
\item Theorem~\ref{thm:local-regularity} (interior regularity) is used in Parts~II--X, and Lemma~\ref{lem:avg} (properties of the averages, among them $\bQ\,d=\dl Q_0$ and the lattice symmetry) in Parts~II--IV, VII, IX and~XI.
\item The decay calculus and the Combes--Thomas lemmas (Lemmas~\ref{lem:decay-calculus}, \ref{lem:CT-bounded} and~\ref{lem:CT-form}) are used in Parts~II, III, VII, VIII, X and~XI.
\item Proposition~\ref{prop:eff-action-gauge}, the comparison of $\bD$ with $\|\dl B\|^2$, is used in Parts~VII and~XI.
\item The operator theorems (Theorems~\ref{thm:scalar-green}, \ref{thm:gauge-coercive}, \ref{thm:gauge-green}, \ref{thm:H-gauge}, \ref{thm:C-gauge}, Propositions~\ref{prop:QGQ}, \ref{prop:R}, \ref{prop:QGQ-gauge}) are the unperturbed operators of Parts~II and~III, and Part~XI uses them directly, in infinite volume (Section~\ref{subsec:ol-infinite-volume}).
\end{itemize}
The constants of Lemmas~\ref{lem:block-poincare}, \ref{lem:biharmonic} and~\ref{lem:Q-divergence} and Corollary~\ref{cor:poincare-all} are tracked in the proof of Lemma~\ref{lem:uniformity}, on the uniformity of the constants of Parts~II and~III. All of Part~I is on the main line.

\paragraph{Details for this Part.} The companion has no chapter on Part~I; its parameter ledger lists the constants of Part~I that Parts~VII--XI use.
